% Hong Kong and Digital Humans ... orientation essay
% Source of truth: places.html (zhao-langxi). UK English. No em dash.
% Compile with tectonic or XeLaTeX (CJK titles need a Unicode engine).
\documentclass[11pt,a4paper]{article}

\usepackage[margin=1in]{geometry}
\usepackage{fontspec}
\usepackage{xeCJK}
\usepackage{setspace}
\usepackage{enumitem}
\usepackage[hidelinks]{hyperref}
\usepackage[british]{babel}

\setmainfont{Times New Roman}
\setCJKmainfont{Songti SC}

\onehalfspacing
\setlength{\parskip}{0.55em}
\setlength{\parindent}{0pt}
\setlist[itemize]{leftmargin=1.4em, itemsep=0.35em, topsep=0.4em}
\Urlmuskip=0mu plus 1mu

\title{Hong Kong and Digital Humans}
\author{Jade Zhao}
\date{September 2026 \\ \small Essay door: \texttt{zhao-langxi} · work in progress · not peer-reviewed}

\begin{document}
\maketitle

\noindent\textit{Essay · Hong Kong focus · not peer-reviewed}

Looking human is not the product. The open question is still trust on an ordinary day: when a system looks and talks like a person, what makes people use it, read the disclosure, and keep using it for real work.
This note looks at Hong Kong as one place where that stack shows up in public catalogues and local demos. It foreshadows study, not a move or a launch.

\section*{Why Hong Kong}

Hong Kong sits between English-facing enterprise tooling and Chinese-language production habits.
Cantonese, Mandarin, and English show up together in local digital-human product pages.
The city's public AI catalogue also frames digital humans as customer-service and enquiry interfaces for departments, not only as entertainment faces.

The Smart Government Innovation LAB (Cyberport) lists solutions such as China Mobile Hong Kong's AI digital-human customer service, Tencent Cloud Digital Human, Momentum Cloud's expressive virtual-human engine (Computer And Technologies), and 01.AI (HK) Limited's avatar / voice-cloning package.
Separately, the Hong Kong Productivity Council's Digital DIY portal lists DYXnet's 線靈 / Digital Human Platform for AI anchors and real-time assistants with Mandarin, Cantonese, and English.

\begin{itemize}
  \item Smart LAB · AI Digital human customer service (China Mobile Hong Kong):
    \url{https://www.smartlab.gov.hk/en/ai_solutions/a-0073}
  \item Smart LAB · Tencent Cloud Digital Human:
    \url{https://www.smartlab.gov.hk/en/ai_solutions/a-0077}
  \item HKPC Digital DIY · DYXnet Digital Human Platform:
    \url{https://ddiy.hkpc.org/en/node/1519}
  \item DYXnet product page (Hong Kong):
    \url{https://www.dyxnet.com/hk/digital-human-platform/}
\end{itemize}

Honest read: much of what is marketed is still broadcast video, virtual anchors, and service-hall Q\&A with a face on top of a knowledge base or large language model.
That is closer to useful business interfaces than pure livestream idols \ldots{} and still not the same as deep work on trust, disclosure, and anthropomorphism.

\section*{Videos worth watching}

Titles and URLs below were checked against YouTube's oEmbed API. They are orientation demos, not endorsements.

\begin{itemize}
  \item DYXnet - Digital Human AI 數字人 video \\
    \url{https://www.youtube.com/watch?v=XVA4cLqBN3g}
  \item Official Cyberport Hong Kong channel · related AIGC public demo (film / creator platform, not a government enquiry avatar):
    【數碼港智慧生活初創培育計劃：創科啟航 · FizzDragon】創意無界：以 AIGC 連結全球創作者 \\
    \url{https://www.youtube.com/watch?v=Kuav2SDlQw4}
  \item Cantonese maker culture around digital humans (consumer tool demo, not an enterprise case study):
    【最逼真1:1 廣東話虛擬人】20秒教識你打造自己一模一樣の虛擬數字人幫你去拍片 \\
    \url{https://www.youtube.com/watch?v=8t9v9TMb-M0}
\end{itemize}

\section*{Mainland China, briefly}

Across the border, mainland 数字人 platforms (for example Baidu's digital-employee / 曦灵 stack and Huawei Cloud MetaStudio) skew harder toward video, livestream, and brand spokespeople, with intelligent interaction as an add-on.
Hong Kong's public catalogue language is more often ``enquiry'', ``customer service'', and multilingual kiosks.
That contrast is useful. It does not make either place a finished answer to trust on an ordinary day.

\section*{Two other hubs for contrast}

\begin{itemize}
  \item \textbf{London · Synthesia} \ldots{} enterprise video moving into interactive avatars for sales, support, and trainers.
    \url{https://jobs.ashbyhq.com/synthesia/e934a7ac-b668-46c0-9953-eed4504658b7}
  \item \textbf{San Francisco / Auckland · Soul Machines} \ldots{} Digital Workforce framed for CX, HR, and related face-to-face workflows.
    \url{https://www.soulmachines.com/digital-workforce}
\end{itemize}

\section*{What this note does not claim}

\begin{itemize}
  \item Not a recommendation to move, claim an internship, or join any named firm.
  \item Not market-size figures or vendor preference percentages.
  \item Not clinical, health-system, or lab-pipeline framing.
  \item Not a claim that every public-service avatar already earns trust.
\end{itemize}

\section*{Related}

The open question this map sits under is on The question:
\url{https://zhao-langxi.github.io/zhao-langxi/question.html}.
Shipped code proof lives on jadexzhao:
\url{https://jadexzhao.github.io/jadexzhao/}.

\vspace{1.2em}
\noindent\small Independent notes \ldots{} not IU-endorsed research and not peer-reviewed publications.\\
Jade Zhao · 赵郎溪 · Digital Humans Project Lead · Expected May 2027

\end{document}
